% Review draft for the quantification/calibration rewrite.
% Approved text is accumulated here before any edit to the main paper.

\section{Calibration}
\label{sec:quant_calibration}

I organize the model parameters into two groups. The first consists of
parameters assigned from external evidence, measured directly in the data, or
fixed as maintained restrictions.
The second consists of parameters estimated jointly by matching model moments
to their empirical counterparts. I describe the two groups in turn, explain
which moments are particularly informative about the internally estimated
parameters, and then report the model's fit.

\subsection{Parameters Set Outside the Estimation}

\paragraph{Demographics and preferences.}
A model period is four years. Households enter adult life at age 18, retire at
age 66, and live at most through age 85. Mortality begins at retirement:
survival is certain before age 66 and thereafter follows four-year
probabilities constructed from the combined male and female survivor counts in
the 2023 Social Security period life table. Dependent children remain in the
household for an expected 18 years, implying a maturity probability of $2/9$
per period. I set risk aversion to $\sigma=2$.

\paragraph{Maintained restrictions.}
I impose that the bequest motive does not vary with family size,
$\theta_n=0$. I also set $\kappa_T=0$, so housing choice is deterministic
conditional on the household state.\footnote{Earlier specifications estimated
$\kappa_T$ jointly with the remaining parameters. Across those searches, the
estimate repeatedly converged to or near zero, while imposing positive values
did not improve model fit. I therefore impose $\kappa_T=0$ in the current
calibration. Future work will revisit this restriction using moments that
directly discipline heterogeneity in tenure and house-size choices.}
